\section*{Bab \ref{ch:b3:genomics} --- Genomika dan Pengurutan}

\begin{solution}{exo:b3:genomics:1}
Sebuah \omterm{def:b3:genomics:sanger}{dideoksinukleotida} tidak memiliki hidroksil 3$'$, sehingga tak ada
ikatan fosfodiester yang dapat dibuat ke nukleotida berikutnya dan rantainya
berhenti. Dengan hanya bentuk dideoksi setiap rantai akan berhenti pada
kedudukan pertama; dengan hanya bentuk normal tak satu pun akan berhenti.
Campurannya menjadikan penghentian itu peristiwa acak pada tiap kedudukan,
sehingga hasilnya membentuk tangga yang lengkap, satu panjang tiap basa
cetakannya.
\end{solution}

\begin{solution}{exo:b3:genomics:2}
$4\times 10^{8}\times \qty{300}{bp} = 1.2\times 10^{11}$ basa. Manusia:
$1.2\times 10^{11}/3.2\times 10^{9} \approx 38\times$. Bakteri:
$1.2\times 10^{11}/5\times 10^{6} = \num{24000}\times$.
\end{solution}

\begin{solution}{exo:b3:genomics:3}
Sebuah kontig adalah urutan bersinambung yang dirakit dari \omterm{def:b3:genomics:ngs}{bacaan} yang
bertindihan; sebuah \omterm{def:b3:genomics:assembly}{perancah} adalah himpunan kontig yang terurut dan
terorientasi dengan celah yang ukurannya ditaksir di antaranya; \omterm{def:b3:genomics:assembly}{N50} adalah
panjang kontig ketika kontig yang telah diurutkan mencapai separuh basa yang
terakit. Di sini kesepuluh kontig \qty{8}{Mb} saja sudah memuat \qty{80}{Mb},
melewati titik tengah \qty{50}{Mb} (kontig ketujuh mencapai \qty{56}{Mb}):
\omterm{def:b3:genomics:assembly}{N50} $= \qty{8}{Mb}$.
\end{solution}

\begin{solution}{exo:b3:genomics:4}
\omterm{prop:b3:genomics:sizes}{Ukuran genom} tidak mengikuti kerumitan organismenya: bawang memiliki lima
kali DNA manusia, ikan paru empat puluh kali; khamir dan cacing berbeda
sepuluh kali lipat DNA-nya dengan cacah gen yang serupa. Kelebihannya bersifat
tak penyandi: unsur yang dapat berpindah beserta sisanya, intron, ulangan
satelit.
\end{solution}

\begin{solution}{exo:b3:genomics:5}
$e^{-c} = 10^{-6}$ memberi $c = 6\ln 10 = 13.8$. Untuk $c = 8$: $N =
cG/L = 8\times 5\times 10^{6}/150 \approx \num{267000}$ \omterm{def:b3:genomics:ngs}{bacaan}, $\theta
= 30/150 = 0.2$, kontignya $= N e^{-6.4} = \num{267000}\times 0.00166
\approx 440$.
\end{solution}

\begin{solution}{exo:b3:genomics:6}
$P(X < 8) = P(X \le 7) \approx P\bigl(Z < (7.5 - 15)/2.74\bigr) =
P(Z < -2.74) \approx 0.003$. Pada $30\times$ kedua alel sebuah
heterozigot hampir selalu terlihat berkali-kali, cakupannya tak rata
(daerah kaya GC memperoleh lebih sedikit \omterm{def:b3:genomics:ngs}{bacaan}, sehingga sebagian tapak
hanya melihat separuh rata-ratanya), dan \omterm{def:b3:genomics:ngs}{bacaan} yang tersisa masih cukup untuk
memisahkan alel yang sungguhan dari galat $10^{-3}$.
\end{solution}

\begin{solution}{exo:b3:genomics:7}
Setiap \omterm{def:b3:genomics:ngs}{bacaan} \qty{150}{bp} dari dalam ulangannya identik dari salinan mana
pun ia berasal, sehingga perakitnya membangun satu simpul \qty{6}{kb} dengan
$500$ lintasan masuk dan $500$ lintasan keluar; ia tak dapat mengetahui
pasangan mana masuk dengan mana keluar, dan \omterm{def:b3:genomics:assembly}{perakitannya} patah menjadi $500$
celah, dengan ulangannya hadir sekali dengan cakupan $500$ kali rata-ratanya.
\omterm{def:b3:genomics:ngs}{Bacaan} yang lebih panjang daripada ulangannya ditambah sisi tunggal pada
kedua ujungnya --- \qty{8}{kb} atau lebih --- melintasi tiap salinan dan
menguraikannya.
\end{solution}

\begin{solution}{exo:b3:genomics:8}
$4.8\times 10^{7}/1000 = \num{48000}$ selisih penyandi; sepertiganya,
sekitar \num{16000}, mengubah proteinnya.
\end{solution}

\begin{solution}{exo:b3:genomics:9}
Harapan cacah positif palsunya $10^{6}\times 5\times 10^{-8} = 0.05$. Risiko
nisbi $1.15$ pada penyakit berfrekuensi \qty{2}{\%} mengubah beberapa kasus
tiap seribu; seribu orang memuat sekitar dua puluh kasus, jauh terlalu sedikit
untuk melihatnya (dayanya tumbuh dengan cacah kasus dan dengan kuadrat
pengaruhnya). Bagi seorang individu, $0.3$ titik persen tak
berarti. Namun alelnya menandai sebuah gen yang perubahan kecil keaktifannya
mengubah risiko penyakit --- gen itu, dan lintasannya, mungkin menjadi
sasaran obat yang penghambatan tuntasnya berpengaruh besar.
\end{solution}

\begin{solution}{exo:b3:genomics:10}
Genetika populasi: pada bakteri, dengan populasi yang raksasa, seleksi
menyingkirkan bahkan sisipan yang sedikit merugikan; pada mamalia, dengan
populasi efektif yang kecil, hanyutan membiarkan DNA tak penyandi yang agak
merugikan menumpuk --- didukung oleh korelasi terbalik antargaris keturunan
antara \omterm{prop:b3:genomics:sizes}{ukuran genom} dan ukuran populasi, dilemahkan oleh kekecualiannya.
Struktur: genom eukariot diserbu transposon yang justru dibersihkan
bakteri, yang cenderung menghapus dan tak memberi perlindungan meiosis bagi
unsur yang mementingkan diri --- didukung oleh korelasi \omterm{prop:b3:genomics:sizes}{ukuran genom}
dengan kandungan transposonnya. Pengaturan: perkembangan yang rumit menuntut
lebih banyak DNA pengatur --- didukung oleh \omterm{def:b3:genomics:comparative}{unsur tak penyandi yang lestari}
di sekitar gen perkembangan, dilemahkan oleh kenyataan bahwa hanya
\qty{8}{\%} genom manusia yang memperlihatkan seleksi sama sekali, sehingga
sebagian besar DNA tak penyandinya tidak bersifat mengatur.
\end{solution}

\begin{solution}{exo:b3:genomics:11}
Awalan kedua jenisnya jatuh dengan kerapatan total $\rho = (N_{1} +
N_{2})/G$. Sebuah \omterm{def:b3:genomics:ngs}{bacaan} sepanjang $L_{i}$ menjadi \omterm{def:b3:genomics:ngs}{bacaan} paling kanan pada
kontignya bila tak ada \omterm{def:b3:genomics:ngs}{bacaan} jenis mana pun yang mulai di dalam $L_{i} - T$
basa sesudahnya: peluangnya $e^{-\rho(L_{i} - T)}$. Jadi kontignya $= N_{1}
e^{-\rho(L_{1} - T)} + N_{2} e^{-\rho(L_{2} - T)}$. Sebuah \omterm{def:b3:genomics:ngs}{bacaan} panjang
menjadi \omterm{def:b3:genomics:ngs}{bacaan} paling kanan dengan peluang yang secara eksponensial lebih
kecil daripada \omterm{def:b3:genomics:ngs}{bacaan} pendek, dan tiap \omterm{def:b3:genomics:ngs}{bacaan} panjang juga menghapus peluang
munculnya celah sepanjang seluruh panjangnya; basa sebanyak itu berupa \omterm{def:b3:genomics:ngs}{bacaan}
pendek menambahkan banyak titik awal tetapi tiap jendela terlindungnya pendek,
sehingga \omterm{def:b3:genomics:ngs}{bacaan} panjanglah yang unggul.
\end{solution}

\begin{solution}{exo:b3:genomics:12}
Dua \omterm{def:b3:genomics:comparative}{penggandaan seluruh genom} meramalkan bahwa pola empat kali lipat itu
berlaku di seluruh genom: kuartet segmen kromosom \omterm{def:b3:genomics:comparative}{paralog} (paralogon) yang
mengusung famili gen yang sama dengan urutan yang sama, dengan penggandaannya
bertarikh sama bagi semua familinya; satu gugus saja pada amfioksus dan
\emph{Ciona}, yang menyimpang sebelum peristiwanya; serta pada lamprei, yang
menyimpang di sekitar peristiwanya, keadaan antara atau keadaan yang berasal
secara bebas. Penggandaan bebas gugus Hox saja meramalkan paralogon hanya
bagi Hox, tetangga yang riwayatnya berbeda, dan tarikh penggandaan yang
berbeda antarfamili. Paralogon seluruh genom dan penanggalan yang sama,
seperti yang teramati, mendukung \omterm{def:b3:genomics:comparative}{penggandaan seluruh genom}.
\end{solution}

\begin{solution}{pb:b3:genomics:1}
\textbf{1.} $0.002\times 6\times 10^{8} = 1.2\times 10^{6}$ \omterm{def:b3:genomics:ngs}{bacaan};
$c = 1.2\times 10^{6}\times 150/4.6\times 10^{6} \approx 39$.
\textbf{2.} $e^{-39} \approx 10^{-17}$: praktis tak ada basa yang tak
terliput.
\textbf{3.} $\theta = 0.2$; kontignya $= 1.2\times 10^{6} e^{-31} \approx
0$: satu kontig secara teori, celahnya hanya pada ulangan.
\textbf{4.} $c = \num{30000}\times 150/4.6\times 10^{6} = 0.98$;
tak terliput $e^{-0.98} = 0.38$, kira-kira \qty{1.7}{Mb}; kontignya $= \num{30000}
\times e^{-0.78} \approx \num{13700}$.
\textbf{5.} Ketujuh operonnya memberi \omterm{def:b3:genomics:ngs}{bacaan} yang identik lalu runtuh menjadi
satu kontig \qty{5}{kb}, yang dimasuki tujuh sisi kiri tunggal dan
ditinggalkan lewat tujuh sisi kanan yang pasangannya tak diketahui: $14$
ujung kontig, tujuh celah.
\textbf{6.} Sekitar \num{4600} gen; penyandinya $0.88\times 4.6\times 10^{6}
= 4.0\times 10^{6}$ bp, gen rata-ratanya \qty{880}{bp}.
\textbf{7.} $0.998\times 6\times 10^{8}\times 150/3.2\times 10^{9}
\approx 28\times$.
\textbf{8.} Sekitar $14$ \omterm{def:b3:genomics:ngs}{bacaan} tiap alel.
\textbf{9.} $28\times 10^{-3}/3 \approx 0.009$ \omterm{def:b3:genomics:ngs}{bacaan} menampilkan suatu
basa keliru tertentu --- peluang tiga atau lebih kira-kira $10^{-7}$ --- dan
\qty{20}{\%} dari $28$ adalah $5.6$ \omterm{def:b3:genomics:ngs}{bacaan}; sebuah galat tak lolos uji
mana pun, sedangkan alel heterozigot yang sungguhan diharapkan pada $14$.
\textbf{10.} $3.2\times 10^{9}/1500 \approx 2.1\times 10^{6}$
tapak heterozigot; pada urutan penyandi $4.8\times 10^{7}/1500 \approx
\num{32000}$.
\textbf{11.} Sebuah \omterm{def:b3:genomics:ngs}{bacaan} \qty{150}{bp} dari sebuah Alu cocok dengan banyak
dari $1.1$ juta salinannya nyaris sama baiknya, sehingga \omterm{def:b3:genomics:ngs}{bacaan} itu
ditempatkan pada salinan yang keliru atau diberi keyakinan pemetaan yang
rendah. Selisih antarsalinannya lalu tampak seperti varian heterozigot, dan
varian yang sungguhan tersembunyi di antaranya: \omterm{met:b3:genomics:pipeline}{pemanggilan varian} di dalam
Alu biasanya disaring keluar, dan unsur itu menjadi titik buta pengurutan
berbacaan pendek.
\textbf{12.} $1.2\times 10^{-8}\times 6.4\times 10^{9} \approx 77$ mutasi
baru. Sekitar $14$ \omterm{def:b3:genomics:ngs}{bacaan} semestinya menampilkan variannya pada anak itu;
tetapi bila tapak seorang tetua hanya diliputi lima \omterm{def:b3:genomics:ngs}{bacaan}, sebuah alel
heterozigot terlewat dengan peluang $2^{-5} = 3\,\%$ dan varian yang diwarisi
salah dipanggil sebagai varian baru --- sehingga tetuanya harus diurutkan
sedalam anaknya.
\textbf{13.} $4.8\times 10^{7}\times 100/150 = 3.2\times 10^{7}$ \omterm{def:b3:genomics:ngs}{bacaan},
dua puluh kali lebih sedikit daripada genomnya: biaya pengurutannya turun
dua puluh kali lipat, lebih besar daripada biaya langkah penangkapannya.
\textbf{14.} $1.1\times 10^{6}\times 300 = 3.3\times 10^{8}$ bp, kira-kira
\qty{10}{\%} genomnya; \qty{10}{\%} \omterm{def:b3:genomics:ngs}{bacaannya}, sekitar $6\times
10^{7}$.
\textbf{15.} $\num{60000}\times 1500 = 9\times 10^{7}$ bp; $9\times
10^{7}/1.6\times 10^{10} = 0.56\,\%$.
\textbf{16.} $0.012\times 2.9\times 10^{9} = 3.5\times 10^{7}$
substitusi, $1.7\times 10^{7}$ tiap garis keturunan; lajunya $1.7\times
10^{7}/(2.9\times 10^{9}\times 6.5\times 10^{6}) = 9\times 10^{-10}$
tiap basa tiap tahun, $2.3\times 10^{-8}$ tiap generasi 25 tahun ---
kira-kira dua kali laju silsilah $1.2\times 10^{-8}$, yakni percanggahan yang
telah dikenal dan menyarankan generasi masa lalu yang lebih panjang atau
perpisahan yang lebih tua.
\textbf{17.} Empat salinan dari \num{16000} akan menjadi \num{64000};
\num{20000} tersisa, sehingga dari \num{48000} salinan tambahannya hanya
\num{4000} yang bertahan: \qty{92}{\%} gandanya hilang.
\textbf{18.} Cacing dan manusia memiliki cacah gen yang sama; kerumitannya
terletak pada cara gen itu dipakai --- penyambungan alternatif mengalikan satu
gen menjadi ribuan protein (\num{38016} bagi \emph{Dscam}), pengaturan
menggabungkan faktor transkripsi dan penguat dengan cara yang khas tipe sel,
\omterm{def:b3:rna-regulation:ncrna}{RNA tak penyandi} menambahkan lapisan, dan proteinnya berinteraksi di dalam
jejaring.
\textbf{19.} Unsur pengatur (penguat, promotor, isolator),
gen \omterm{def:b3:rna-regulation:ncrna}{RNA tak penyandi}, isyarat penyambungan dan penempatan, titik asal, serta
urutan struktural. Ujilah sebuah calon penguat dengan menempatkannya di depan
gen pelapor di dalam embrio transgenik lalu melihat di mana pelapornya
diekspresikan, kemudian menghapus unsurnya di dalam genom dengan CRISPR lalu
mengukur gen tetangganya.
\textbf{20.} $0.05/10^{6} = 5\times 10^{-8}$; pada $p < 10^{-5}$, $10^{6}
\times 10^{-5} = 10$ positif palsu diharapkan.
\textbf{21.} Ods bukan pembawanya $0.009/0.991 = 0.00908$; ods satu salinan
$0.00954$, risikonya $0.945\,\%$; ods dua salinan $0.01001$, risikonya
$0.99\,\%$.
\textbf{22.} Skornya menjumlahkan, atas $200$ alelnya, cacah salinan
risiko yang diusung seseorang dengan bobot log rasio ods tiap alel. Cacahnya
berrata-rata $200\times 2\times 0.3 = 120$ dan bersimpangan baku
$\sqrt{200\times 2\times 0.3\times 0.7} \approx 9$; \qty{1}{\%} teratasnya
mengusung sekitar $21$ alel lebih banyak daripada rata-rata, dan $1.05^{21} \approx
2.8$: nyaris tiga kali lipat ods orang rata-rata, dari alel
yang secara sendiri-sendiri mengubah risiko sebesar \qty{5}{\%}.
\textbf{23.} Alel di dalam satu blok sepanjang puluhan kilobasa diwarisi
bersama (ketakseimbangan pautan), sehingga alel mana pun di antaranya
memperlihatkan asosiasi yang sama; SNP yang digenotipe sekadar SNP yang
kebetulan ada pada lariknya. Untuk menemukan varian penyebabnya: urutkanlah
daerah itu pada kasus dan kendalinya, sempitkanlah himpunannya secara
statistik, lalu ujilah tiap calonnya --- uji pelapor bagi keaktifan penguat
tiap alelnya, penyuntingan tiap alel di dalam sel lalu pengukuran ekspresi
gen di dekatnya, penempatan bersama dengan isyarat sifat kuantitatif
ekspresi.
\textbf{24.} $1 - e^{-0.5} = 39\,\%$ genomnya terbaca. Genom itu memberi
contoh langsung sebuah populasi pada waktu yang diketahui di masa lalu ---
frekuensi alel sebelum perpindahan dan pembauran berikutnya, panjang segmen
Neanderthal (panjang, karena baru sedikit generasi rekombinasi yang
memecahnya) dan karenanya tarikh pembaurannya --- yang hanya dapat disimpulkan
genom modern lewat model.
\textbf{25.} Sekitar \num{13700} kontig pada rintisannya; $28\times$ tiap
genom haploid, sekitar $14$ \omterm{def:b3:genomics:ngs}{bacaan} tiap alel; sekitar $77$ mutasi baru pada
seorang anak.
\end{solution}
